\chapter{Obiectivele proiectului}\label{ch:obiective}
\pagestyle{fancy}

\section{Introducere}
Acest capitol prezintă obiectivele proiectului. Ele sunt legate direct de problema adnotării vizuale și descriu ce se urmărește prin realizarea aplicației. Structura include obiectivul general, obiectivele principale și secundare, precum și cerințele funcționale și non-funcționale.

\section{Obiectivul general}
În mod obișnuit, adnotarea datelor vizuale este o muncă manuală, costisitoare și lentă, iar rezultatele pot fi influențate de oboseală și erori umane. Acest context creează nevoia unei soluții care să accelereze procesul, fără a pierde controlul asupra calității.

Obiectivul general al proiectului este dezvoltarea unei aplicații desktop de adnotare vizuală, semi-automată și colaborativă, care reduce timpul de etichetare și păstrează controlul utilizatorului asupra rezultatului final. Aplicația trebuie să ofere un flux de lucru complet, de la încărcarea imaginilor și navigarea rapidă între ele, până la editarea adnotărilor și exportul în formate standard. În mod explicit, trebuie să suporte adnotări de tip bounding box (dreptunghi), poligon și mască (image segmentation), pentru a acoperi scenariile uzuale de detecție și segmentare.

Gestionarea adnotărilor trebuie să fie clară și previzibilă, cu posibilitatea de selecție, editare și ștergere, astfel încât utilizatorul să poată corecta rapid erorile. Rezultatele trebuie să poată fi salvate și exportate în formate standard, pentru a fi folosite ușor în antrenarea modelelor. Prin integrarea unor modele externe pentru propuneri de segmentare și prin suportul pentru colaborare securizată între utilizatori, soluția urmărește să crească eficiența fără a compromite calitatea datelor.

\section{Obiective principale}
Obiectivele principale urmăresc să definească funcționalitățile cheie ale aplicației și modul în care acestea răspund cerințelor de adnotare. Fiecare obiectiv este formulat astfel încât să poată fi implementat și verificat concret.

\begin{itemize}
	\item \textbf{Interfață de adnotare completă.} Realizarea unui canvas interactiv cu zoom și pan, care permite desenarea de adnotări de tip bounding box, poligon și mască. Interfața trebuie să susțină navigarea rapidă între imagini și un control vizual clar.
	\item \textbf{Gestionarea coerentă a adnotărilor.} Păstrarea unei structuri clare pentru etichete și elemente, cu identificatori unici, listare și operații de editare sau ștergere. Această structură reduce erorile și facilitează modificările ulterioare.
	\item \textbf{Export standardizat al datelor.} Oferirea de export în formatele COCO și YOLO, cu organizare predictibilă a fișierelor. Astfel, dataset-urile pot fi folosite direct în antrenare, fără conversii suplimentare.
	\item \textbf{Integrare AI semi-automată.} Folosirea AutoSeg și AutoMask pentru a genera propuneri de segmentare care pot fi corectate imediat de utilizator. Scopul este accelerarea procesului fără a compromite calitatea adnotărilor.
	\item \textbf{Modularizarea execuției AI.} Rularea modelelor ca aplicații separate, prin procese dedicate, pentru a izola erorile și a permite actualizări independente. Această abordare crește stabilitatea și flexibilitatea sistemului.
	\item \textbf{Colaborare în timp real.} Realizarea unui server separat pentru colaborare, cu suport de conectare și sesiuni comune, apoi sincronizare a evenimentelor de tip create, update și delete. Actualizarea imaginii între utilizatori se face prin REST și WebSocket pentru a evita duplicarea muncii.
	\item \textbf{Comunicare între utilizatori.} Oferirea unei interfețe de chat care permite schimbul de mesaje între utilizatori în timpul sesiunilor de lucru.
	\item \textbf{Securitate și acces controlat.} Autentificarea pe bază de JWT și păstrarea locală a token-ului pentru o experiență de lucru fluidă. Accesul la sesiuni și date trebuie să fie controlat și verificabil.
	\item \textbf{Evaluarea și selecția modelelor AI externe.} Testarea Mask2Former și YOLOE pe imagini de referință, cu măsurarea timpului de procesare și a calității rezultatelor. Alegerea finală trebuie justificată pe baza acestor rezultate.
	\item \textbf{Impachetarea și integrarea backend-urilor de inferență.} Folosirea unui script care facilitează build-uirea executabilului, pentru un export simplu și repetabil. Integrarea trebuie să permită conectarea ușoară a backend-urilor la aplicație.
\end{itemize}

\section{Obiective secundare}
Obiectivele secundare completează funcționalitățile principale prin îmbunătățirea experienței de utilizare și prin validări suplimentare ale componentelor AI și de colaborare.

\begin{itemize}
	\item \textbf{Benchmarking pentru modele alternative.} Analizarea unor modele open-vocabulary (YOLO-World, LLMDet, SAM2) pentru a observa potențialul de generalizare și a justifica alegerea finală.
	\item \textbf{Validarea suportului local de tip VLM/LLM.} Compararea mai multor modele rulate local (Ollama) pentru a susține componenta de chatbot și a stabili un compromis calitate-viteză. Chatbotul are rolul de a ajuta utilizatorul cu întrebări și de a oferi sugestii pentru remedieri ale erorilor întâlnite.
	\item \textbf{Experiență de utilizare și prevenirea erorilor.} Confirmare la schimbarea imaginii și la închidere pentru a evita pierderea adnotărilor.
	\item \textbf{Persistența setărilor vizuale.} Salvarea preferințelor (grosime contur, dimensiune font, opacitate mască) între sesiuni.
	\item \textbf{Îmbunătățirea vizuală a imaginii.} Ajustări de luminozitate, contrast și gamma fără a altera datele originale.
	\item \textbf{Încărcare inteligentă a adnotărilor.} Asocierea automată a fișierelor JSON cu imaginea corespunzătoare.
	\item \textbf{Optimizarea serializării măștilor.} Compresie și reducerea volumului de date pentru fișiere mari.
	\item \textbf{Reducerea traficului în colaborare.} Împărțirea actualizărilor mari (ex. măști) în batch-uri și trimiterea lor prin WebSocket, pentru a evita închiderea prematură a conexiunii. Se aplică și deduplicarea evenimentelor.
	\item \textbf{Crosshair și ghidaj vizual.} Indicatori pentru precizie în adnotare și selectare.
\end{itemize}

\newpage

\section{Cerințe funcționale și non-funcționale}
Pentru a susține obiectivele prezentate, proiectul are un set de cerințe funcționale și non-funcționale, sintetizate în tabelele de mai jos.

\subsection{Cerințe funcționale}
\renewcommand{\arraystretch}{1.2}
\begin{center}
\begin{tabular}{|p{0.32\linewidth}|p{0.6\linewidth}|}
\hline
\textbf{Cerință} & \textbf{Descriere scurtă} \\
\hline
Încărcare și navigare imagini & Deschidere folder, listare rapidă și comutare între imagini. \\
\hline
Canvas de adnotare & Zoom, pan și trasare de bounding box, poligon și mască. \\
\hline
Editare adnotări & Selectare, modificare, ștergere, etichete și culori ușor de gestionat. \\
\hline
Import/Export & Salvare JSON și export COCO/YOLO cu structură predictibilă. \\
\hline
AI semi-automată & AutoSeg/AutoMask pentru propuneri rapide, corectabile imediat. \\
\hline
Modularizare AI & Rulare modele ca procese separate pentru izolare și actualizări. \\
\hline
Colaborare & Sesiuni, sincronizare create/update/delete prin REST și WebSocket. \\
\hline
Autentificare & Acces controlat cu JWT și token salvat local. \\
\hline
Comunicare & Chat între utilizatori și chatbot local pentru întrebări și erori. \\
\hline
Personalizare & Setări vizuale persistente și confirmare la închidere/schimbare. \\
\hline
\end{tabular}
\end{center}

\subsection{Cerințe non-funcționale}
\begin{center}
\begin{tabular}{|p{0.32\linewidth}|p{0.6\linewidth}|}
\hline
\textbf{Cerință} & \textbf{Descriere scurtă} \\
\hline
Usabilitate & Interfață clară, operații frecvente accesibile rapid. \\
\hline
Performanță & Interfață responsivă în timpul procesării AI (asincron). \\
\hline
Fiabilitate & Validări pentru căi inexistente și mesaje clare de eroare. \\
\hline
Securitate & Protecția sesiunilor și a datelor prin token. \\
\hline
Integritate date & Exporturi consistente, fără pierderi de informație. \\
\hline
Scalabilitate & Mai mulți utilizatori pot lucra în aceeași sesiune. \\
\hline
Mentenabilitate & Arhitectură modulară pentru întreținere și extindere. \\
\hline
Extensibilitate & Adăugare de modele și formate noi fără refactor major. \\
\hline
Portabilitate & Rulare locală pe Windows cu configurare minimă. \\
\hline
\end{tabular}
\end{center}
\renewcommand{\arraystretch}{1.0}

